% =====================================================================
%  2. OBJETIVOS
%  Responsable: TODOS
%  ESTADO: REDACTADO (1 de agosto de 2026)
%
%  REGLA: cada objetivo específico tiene DESPUÉS su conclusión
%  correspondiente en la sección 9, en el mismo orden. Son pares.
%  Si añades un objetivo, añades una conclusión.
% =====================================================================

\section{Objetivos}
\label{sec:objetivos}

% ---------------------------------------------------------------------
\subsection{Objetivo general}

Analizar el proceso y los resultados de la auditoría socioambiental aplicada a la
Mina Marlin (Montana Exploradora de Guatemala, S.\,A. / Goldcorp), departamento de
San Marcos, Guatemala, con el fin de formular una propuesta de plan de mejora
continua fundamentada en el ciclo PHVA y en los estándares internacionales de
desempeño ambiental y social.

% ---------------------------------------------------------------------
\subsection{Objetivos específicos}

\begin{enumerate}
  \item Describir la unidad minera y su entorno físico, geológico y social,
        identificando los aspectos de la operación que determinan su perfil de
        riesgo socioambiental.
  \item Identificar el alcance, los criterios de auditoría y la metodología
        aplicados por el equipo auditor, y contrastarlos con las directrices de la
        ISO 19011:2018.
  \item Clasificar los hallazgos ambientales y sociales de la auditoría según las
        categorías de no conformidad mayor, no conformidad menor y observación,
        estableciendo para cada uno su evidencia objetiva, el criterio incumplido
        y su estado de cierre.
  \item Evaluar críticamente la independencia del proceso de auditoría, sus
        limitaciones metodológicas declaradas y su eficacia real sobre la conducta
        de la organización auditada.
  \item Proponer un plan de mejora continua estructurado sobre el ciclo PHVA, con
        acciones verificables mapeadas a requisitos concretos de los IFC
        Performance Standards, el GISTM y el estándar IRMA.
\end{enumerate}
